% Part: incompleteness
% Chapter: arithmetization-syntax
% Section: introduction

\documentclass[../../../include/open-logic-section]{subfiles}

\begin{document}

\olfileid[tr]{inc}{art}{int}
\olsection{Giriş}

Hesaplanabilirlik ile mantık arasında bağlantı kurmak için mantığın
nesnelerinden (simgelerden, terimlerden, !!{formula}s{}den,
!!{derivation}s{}den), bunlar üzerindeki işlemlerden, bunların özellik ve
bağıntılarından hesaplamalı işleme elverişli bir biçimde söz edebilmemiz
gerekir. Bunu, simgeler, simge dizileri ve bunlardan kurulan diğer nesneler
üzerindeki hesaplanabilir fonksiyon ve bağıntıları ele alarak doğrudan
yapabiliriz. Mantıksal sözdizimin bütün nesneleri sonlu olduğundan ve
!!a{enumerable} simge kümelerinden kurulduğundan, bazı hesaplama modelleri
için bu mümkündür. Fakat özyinelemeli fonksiyonlar gibi diğer hesaplama
modelleri yalnızca sayılarla ve bunlar üzerindeki bağıntı ve fonksiyonlarla
sınırlıdır. Üstelik en sonunda sözdizimiyle belirli teorilerin, özellikle
aritmetik dilinde formüle edilen teorilerin içinde de uğraşabilmek isteriz.
Bu durumlarda sözdizimi \emph{aritmetikleştirmek}, yani sözdizimsel nesneleri,
bunlar üzerindeki işlemleri ve bunların bağıntılarını sırasıyla sayılar,
aritmetik fonksiyonlar ve aritmetik bağıntılar olarak temsil etmek gerekir.
Leibniz'e kadar uzanan düşünce, sözdizimsel nesnelere sayılar atamaktır.

Simgelere ``kodları'' olarak sayılar atamak görece kolaydır. Örneğin sonsuz
sayıda !!{variable} ve hatta her $n$ aritesi için sonsuz sayıda
!!{function} bulunduğundan, bazı simgeler küçük bir güçlük çıkarır. Ama
elbette simgelere sistemli biçimde, örneğin $\Obj v_2$ ile $\Obj v_3$'e
farklı kodlar atanacak biçimde sayılar vermek mümkündür. Terimler ve
!!{formula}s gibi simge dizileri daha büyük bir güçlüktür. Ne var ki sayı
dizilerini salt aritmetik yoldan (örneğin dizilerin asal sayıların kuvvetleriyle
kodlanması yoluyla) ele alabilirsek, tek tek simgelerin kodlanmasını simge
dizilerinin kodlanmasına, ardından !!{formula}s{}den oluşan dizilerin ya da
!!{derivation}s gibi başka düzenlemelerin kodlanmasına genişletebiliriz.
Bu genişletilmiş kodlamaya ``Gödel numaralandırması'' denir. Her terime,
!!{formula} ve !!{derivation} için bir Gödel numarası atanır.

Simge dizilerini simgelerinin kodlarından oluşan diziler olarak kodlayıp
dizileri hesaplanabilir fonksiyonlarla işlenebilen bir kodlama sistemi
seçtiğimizde, Gödel numaralarını da hesaplanabilir fonksiyonlarla ele
alabiliriz. Uygulamada ilgili bütün fonksiyonlar ilkel özyinelemeli olacaktır.
Örneğin bir dizinin uzunluğunu hesaplamak ve dizinin kodundan $i$'inci
elemanını hesaplamak ilkel özyinelemelidir. Diziyi kodlayan sayı, diyelim ki
!!a{formula} $!A$ için Gödel numarasıysa, bu !!a{formula} için uzunluğun ve
!!a{formula} içindeki $i$'inci simgenin kodunun da $!A$'nın Gödel numarasından
hesaplanabildiğini hemen görürüz. Örneğin bir sayının düzgün kurulmuş bir
terimin ya da doğru bir !!{derivation}{}in Gödel numarası olması özelliğinin
ilkel özyinelemeli olduğunu ispatlamak biraz daha zordur. Bununla birlikte,
ilgili diziler (terimler, !!{formula}s ve !!{derivation}s) tümevarımlı olarak
tanımlandığından bu mümkündür.

Örnek olarak yerine koyma işlemini ele alalım. $!A$ bir formül, $x$ bir
değişken ve $t$ bir terimse $\Subst{!A}{t}{x}$, $!A$ içindeki her serbest
$x$ geçişinin $t$ ile değiştirilmesinin sonucudur. Şimdi $!A$, $x$ ve $t$'ye
sırasıyla $k$, $l$ ve $m$ Gödel numaralarını atadığımızı varsayalım. Aynı
şema $\Subst{!A}{t}{x}$'e de, diyelim $n$ Gödel numarasını atar. $k$, $l$
ve $m$'yi $n$'ye götüren bu eşleme, yerine koyma işleminin aritmetik
karşılığıdır. Yerine koyma işlemi $!A$, $x$, $t$'yi
$\Subst{!A}{t}{x}$'e götürürken aritmetikleştirilmiş yerine koyma fonksiyonu
$k$, $l$, $m$ Gödel numaralarını $n$ Gödel numarasına götürür. Bu
fonksiyonun ilkel özyinelemeli olduğunu göreceğiz.

Sözdizimin aritmetikleştirilmesi yalnızca soyut bakımdan ilginç değildir.
Gerçi aritmetik dili gibi, dillerden ``söz etme'' düzeneklerini kendiliğinden
içermeyen dillerin ifadelerin karmaşık özelliklerini biçimselleştirebildiğini
görmek başlangıçta hiç de önemsiz olmayan bir kavrayıştı. Bundan sonra bu
dildeki bir teorinin, örneğin Peano aritmetiğinin, kendi dili hakkında
(örneğin !!{sentence}s{}in ispatlanabilir ya da doğru olup olmadığı hakkında)
neyi \emph{ispatlayabildiğini} sormak küçük bir adımdır. Bu bizi Gödel'in
(ispatlanamazlık hakkındaki) ve Tarski'nin (doğrunun tanımlanamazlığı
hakkındaki) ünlü sınırlama teoremlerine götürür. Ancak sözdizimi
aritmetikleştirme yöntemi, hesaplanabilirlik teorisinde teorilerin hesaplama
gücü ya da farklı hesaplanabilirlik modelleri arasındaki ilişki gibi önemli
sonuçları ispatlamak için de önemlidir. Sözdizimin aritmetikleştirilmesi,
başka nesne ve özellikleri aritmetikleştirmeye örnek olur. Örneğin
konfigürasyonları ve hesaplamaları (diyelim Turing makinelerininkileri) benzer
biçimde aritmetikleştirmek mümkündür. Bu, bir modeldeki (örneğin Turing
makinelerindeki) hesaplamaların başka bir modelde (örneğin özyinelemeli
fonksiyonlarda) benzetilmesini sağlar.

\end{document}
